% Problem record op_39721600d81c6992. This TeX file is derived from
% database/problems_json/op_39721600d81c6992.json by scripts/sync-tex.mjs; edit the JSON record
% and rerun the script rather than editing this file.
\section{Quantum query complexity of Graph Collision}
\label{sec:op_39721600d81c6992}

\paragraph{Problem.}
Does Graph Collision on every known $n$-vertex graph admit a bounded-error
quantum algorithm using $\widetilde O(\sqrt n)$ queries to the marked-vertex bits?
Let $G=([n],E)$ be a simple undirected graph given in full, and let
$x\in\{0,1\}^n$ be accessible through the quantum oracle
$O_x|v,b\rangle=|v,b\oplus x_v\rangle$.
The task is to compute
\begin{equation}
  f_G(x)=\bigvee_{\{u,v\}\in E}(x_u\wedge x_v).
  \label{eq:3972-collision}
\end{equation}
Only calls to $O_x$ are counted.
All computation involving the known graph
is free, and the algorithm must return the value in
\eqref{eq:3972-collision} with probability at least $2/3$ for every $x$.

Let $Q_2(f_G)$ be the minimum worst-case query count under this guarantee,
and let $Q_{\mathrm{GC}}(n)=\max_{G:\,|V(G)|=n}Q_2(f_G)$.
The question is whether
\begin{equation}
  Q_{\mathrm{GC}}(n)=\widetilde O(\sqrt n).
  \label{eq:3972-target}
\end{equation}
In \eqref{eq:3972-target}, $\widetilde O$ hides only factors polynomial in
$\log n$ as $n\to\infty$.

\subsection*{Status}

Unsolved

\subsection*{Source}

Gavinsky and Ito explicitly ask for a $\widetilde O(\sqrt n)$ query algorithm
for all graphs in Section~6 of \sourcecite{ref:3972-gi}{GI12}.
Their Definition~1 fixes the known-graph, marked-vertex oracle model.

\subsection*{Progress}

\begin{itemize}
  \item Gavinsky and Ito give an
$O(\sqrt n+\sqrt{\alpha^*(G)})$ query algorithm, where
$\alpha^*(G)$ is the maximum sum of vertex degrees over an independent set
(Theorem~1) \sourcecite{ref:3972-gi}{GI12}.
Their Corollary~2 gives $O(\sqrt{n\log n})$ queries with high probability
over a random graph with independently sampled edges, for any choice of edge probability.
This result for typical graphs does not give the target bound for every graph.

  \item Cornelissen, Gilani, Kintali, Le Gall, Patro, and Zhou report an algorithm
using $\widetilde O(n^{11/17})$ queries for arbitrary graphs and a lower bound
of $\Omega(\sqrt{n\log n})$ for a family of graphs
(Theorems~1.1 and~6.1, preprint version~1)
\sourcecite{ref:3972-cgklpz}{CGKLPZ26}.
Both bounds concern queries to marked-vertex bits, with the graph known for free.
The lower bound excludes $O(\sqrt n)$ queries in general, but is compatible
with the polylogarithmic factors allowed in \eqref{eq:3972-target}.
\end{itemize}

\subsection*{References}

\begin{enumerate}[label={},leftmargin=4.8em,labelsep=0.5em]
  \footnotesize
  \item[\textup{[GI12]}]\label{ref:3972-gi}
  Dmitry Gavinsky and Tsuyoshi Ito,
``A quantum query algorithm for the graph collision problem,'' preprint (2012).
\href{https://arxiv.org/abs/1204.1527v1}{arXiv:1204.1527v1}.

  \item[\textup{[CGKLPZ26]}]\label{ref:3972-cgklpz}
  Arjan Cornelissen, Amin Shiraz Gilani, Shiva Kintali, Fran\c{c}ois Le Gall,
Subhasree Patro, and Xingyu Zhou,
``Breaking the quantum barriers for Graph Collision,'' preprint (2026).
\href{https://arxiv.org/abs/2611.05356v1}{arXiv:2611.05356v1}.
\end{enumerate}

\subsection*{Comment}

The remaining question is whether the upper bound can be reduced to
$\widetilde O(\sqrt n)$, or whether some graph family requires more than
$\sqrt n$ times every fixed power of $\log n$ queries.
Graph Collision here is a query problem with a known graph.
It is distinct from triangle finding with oracle access to unknown edges,
and its query bound does not assert a matching running time.

\subsection*{Field}

Quantum algorithm

\subsection*{Topic}

Computational complexity and computability

\subsection*{ID}

\texttt{op\_39721600d81c6992}
